\clearpage\pagestyle{otMain}\setlength{\otBlocEspacement}{2.5mm}\fontsize{11}{15}\selectfont\corrigesChapitre{12}{Les aires : unités et formules}
\begin{blocLivre}{cor-12-1}\begin{corrige}{1}{Compter puis calculer}\demarche{Relier rangées et produit.}
Comptage : 4 rangées de 7, soit $\res{28}$ carreaux. Formule : $7\times4=\res{28\ \mathrm{cm}^2}$ : la formule résume le pavage.
\label{fig-cor-12-1}\figureGeo[corrige]{aire-1}
\end{corrige}\end{blocLivre}
\begin{blocLivre}{cor-12-2}\begin{corrige}{2}{Carré et rectangle}\demarche{Multiplier les deux dimensions.}
Carré : $6\times6=\res{36\ \mathrm{cm}^2}$. Rectangle : $9\times4{,}5=\res{40{,}5\ \mathrm{cm}^2}$.
\end{corrige}\end{blocLivre}
\begin{blocLivre}{cor-12-3}\begin{corrige}{3}{Aire ou périmètre ?}\demarche{Associer chaque formule à sa grandeur.}
Périmètre : $4\times6=\res{24}$ cm, une longueur. Aire : $6\times6=\res{36}\ \mathrm{cm}^2$, une surface. Deux grandeurs, deux unités.
\end{corrige}\end{blocLivre}
\begin{blocLivre}{cor-12-4}\begin{corrige}{4}{Les conversions}\demarche{Multiplier ou diviser par 100.}
Le rapport est 100 : $\res{300\ \mathrm{cm}^2}$ ; $\res{450\ \mathrm{dm}^2}$ ; $\res{2\ \mathrm{dm}^2}$.
\end{corrige}\end{blocLivre}
\begin{blocLivre}{cor-12-5}\begin{corrige}{5}{Pourquoi 100 ?}\demarche{Paver le grand carré par les petits.}
Le carré de 1 dm de côté contient $\res{10}$ rangées de $\res{10}$ carreaux de 1 cm : $10\times10=\res{100}$ carreaux. L’aire croise deux longueurs : le rapport 10 se multiplie par lui-même.
\label{fig-cor-12-5}\figureGeo[corrige]{aire-5}
\end{corrige}\end{blocLivre}
\begin{blocLivre}{cor-12-6}\begin{corrige}{6}{Retrouver un côté}\demarche{Défaire la multiplication.}
Largeur : $54:9=\res6$ cm. Pour le carré : $7\times7=49$, donc côté $\res7$ cm.
\end{corrige}\end{blocLivre}
\begin{blocLivre}{cor-12-7}\begin{corrige}{7}{La figure en L}\demarche{Compléter ou découper, puis comparer.}
Par différence : $8\times6-5\times4=48-20=\res{28\ \mathrm m^2}$. Par découpage : un rectangle $8\times2=16$ et un rectangle $3\times4=12$ : $16+12=28\ \mathrm m^2$ aussi.
\label{fig-cor-12-7}\figureGeo[corrige]{aire-7}
\end{corrige}\end{blocLivre}
\begin{blocLivre}{cor-12-8}\begin{corrige}{8}{Les unités d’abord}\demarche{Convertir avant la formule.}
En centimètres : $200\times50=\res{10\,000\ \mathrm{cm}^2}$. Or $1\ \mathrm m^2=10\,000\ \mathrm{cm}^2$ : l’aire vaut $\res{1\ \mathrm m^2}$. La demi-longueur compense le double mètre.
\end{corrige}\end{blocLivre}
\begin{blocLivre}{cor-12-9}\begin{corrige}{9}{Même périmètre, aires rangées}\demarche{Calculer les deux grandeurs séparément.}
Périmètres : $2\times(9+1)=2\times(7+3)=2\times(5+5)=20$ cm. Aires : $9$, $21$ et $25\ \mathrm{cm}^2$ : $\res{9<21<25}$. À périmètre fixé, le carré offre la plus grande aire.
\end{corrige}\end{blocLivre}
\begin{blocLivre}{cor-12-10}\begin{corrige}{10}{Le carré agrandi}\demarche{Comparer les deux aires par leur quotient.}
Avant : $4\times4=16\ \mathrm{cm}^2$. Après : $12\times12=144\ \mathrm{cm}^2$. Rapport : $144:16=\res9$ : tripler le côté multiplie l’aire par $3\times3=9$.
\end{corrige}\end{blocLivre}
\begin{blocLivre}{cor-12-11}\begin{corrige}{11}{La chambre à moquette}\demarche{Calculer l’aire avant le prix.}
Aire : $4{,}2\times3{,}5=\res{14{,}7\ \mathrm m^2}$. Prix : $14{,}7\times18=\res{264{,}60}$ euros.
\end{corrige}\end{blocLivre}
\begin{blocLivre}{cor-12-12}\begin{corrige}{12}{Le potager et l’allée}\demarche{Soustraire la partie non cultivée.}
Terrain : $12\times12=144\ \mathrm m^2$. Allée : $12\times2=24\ \mathrm m^2$. Cultivé : $144-24=\res{120\ \mathrm m^2}$.
\end{corrige}\end{blocLivre}
\begin{blocLivre}{cor-12-13}\begin{corrige}{13}{Les dalles de la terrasse}\demarche{Compter les dalles par rangées.}
Une dalle couvre $50\times50=2\,500\ \mathrm{cm}^2$, soit $0{,}25\ \mathrm m^2$. Terrasse : $24\ \mathrm m^2$. Nombre : $24:0{,}25=\res{96}$ dalles ; en effet $12\times8=96$ dalles en comptant 12 par longueur et 8 par largeur.
\end{corrige}\end{blocLivre}
\begin{blocLivre}{cor-12-14}\begin{corrige}{14}{Deux annonces}\demarche{Mettre les deux surfaces en m².}
B : $9\times5{,}2=46{,}8\ \mathrm m^2$. B dépasse A de $46{,}8-45=\res{1{,}8\ \mathrm m^2}$.
\end{corrige}\end{blocLivre}
\begin{blocLivre}{cor-12-15}\begin{corrige}{15}{Vrai ou faux ?}\demarche{Construire un contre-exemple chiffré.}
Aires : $12\times3=36$ et $9\times4=36$ : égales. Périmètres : $2\times15=30$ et $2\times13=26$ : différents. L’affirmation est $\res{\text{fausse}}$.
\end{corrige}\end{blocLivre}
